% id: 66-37
% book: 베이직쎈 확률과 통계 · 04 조건부확률
% section: 개념 19 $\mathrm{P}(B)=\mathrm{P}(A\cap B)+\mathrm{P}(\comp{A}\cap B)$를 이용한 확률
% passage: 다음에 답하시오.
% figure: none
% answer: 
% answer_source: 
% variant_level: 0 (원본 전사)
% 이 파일은 build.mjs 가 items.json 에서 만든다. 고칠 때는 items.json 을 고친다.
\smash{\raisebox{1.5mm}{\dmkichul{베이직쎈 확률과 통계 66쪽 37번}}}
꿀이 들어 있는 송편 2개를 포함하여 7개의 송편이 담겨 있는 $\mathrm{A}$ 접시와 꿀이 들어 있는 송편 2개를 포함하여 6개의 송편이 담겨 있는 $\mathrm{B}$ 접시가 있다. 두 접시 $\mathrm{A}$, $\mathrm{B}$ 중 임의로 하나를 택하여 한 개의 송편을 먹었더니 꿀이 들어 있는 송편이었을 때, 택한 접시가 $\mathrm{B}$일 확률을 구하시오.
